\documentclass[10pt]{article}
\usepackage[margin=0.7in]{geometry}
\usepackage[T1]{fontenc}
\usepackage{lmodern}
\usepackage{xcolor}
\usepackage{hyperref}
\usepackage{amsmath,amssymb}

\begin{document}

\title{\vspace{-2.0cm}Response to Reviewer 1}
\author{Christopher Waight and Christopher A. Kitts}
\date{\today}
\maketitle

\vspace{-0.5cm}
We thank the reviewer for their thorough and constructive feedback. We appreciate the detailed review of the mathematics, benchmark analysis, and claims. Below is our point-by-point response addressing every single piece of feedback in the report, including minimal textual fixes proposed for the next revision.

\vspace{-0.2cm}
\section*{2.2 Where the mathematics can be challenged}

\textbf{M1. The rotating-observer model matches one sensing pipeline, and the paper does not say which} \\
\textit{Response:} We initially questioned whether the $\mathbf{v}' = \mathbf{Q}\mathbf{v} + \Omega R \mathbf{p}'$ kinematic artifact accurately applied, but we appreciate the rigorous breakdown demonstrating its exact correspondence. The mathematical derivation provided by the reviewer is sound for the specific sensing pipeline where a robot measures relative water velocity and adds its own velocity obtained from position fixes in the drifting axes. \\
\textbf{Minimal Fix Proposed:} 
\begin{enumerate}
    \setlength{\itemsep}{0pt}
    \setlength{\parskip}{0pt}
    \item Add two sentences to IV-D right after the rotating-frame field $\mathbf{v}'$: \textit{``This is the field a cluster measures when each robot senses water velocity relative to itself and adds its own velocity from position fixes in the drifting axes. A sensor that reports an earth-referenced vector rotated by the heading error measures $\mathbf{Q}\mathbf{v}$ instead, which leaves $D$ unchanged.''}
    \item In the Conclusion, narrow the rule: \textit{``One whose flow estimate is referenced to a drifting frame...''}
    \item Replace ``about $7^\circ$/h'' with the dimensionless rate $\Omega/(\pi^2 A) \approx 0.21$.
\end{enumerate}

\vspace{-0.2cm}
\section*{Stability, Heuristics, and Noise}

\textbf{M2. The stability argument is a continuous-time idealization} \\
\textit{Response:} We agree with the scope regarding the positive constant $a_\perp$. \\
\textbf{Minimal Fix Proposed:} Add the scoping sentence to the Appendix before the Proposition: \textit{``The analysis is for the continuous-time idealization; where $a_\perp$ varies along the ride, the bound holds with $a_\perp$ replaced by its infimum.''} Add the uniform-decrease step to the proof, and the one-sentence benchmark profile to III-B.

\textbf{M3. The $D$ tracker's capture test cannot certify its own target on the benchmark} \\
\textit{Response:} We agree the fitted Hessian lacks the terms to make it definite at a saddle. \\
\textbf{Minimal Fix Proposed:} After Eq. (18), add: \textit{``The fitted Hessian lacks the $\mathbf{J}\cdot\partial^3\mathbf{v}$ terms that make $\mathbf{H}_D$ definite at a saddle, so (18) is a heuristic.''}

\textbf{M4. The $s_1$ noise fragility has a per-cycle explanation that the paper sets aside} \\
\textit{Response:} The reviewer rightly separates persistence from accuracy. \\
\textbf{Minimal Fix Proposed:} Soften the attribution in IV-C to: \textit{``Noise flips the sign at tangent selection, and because the $s_1$ tracker carries that sign in state through $\mathbf{t}_{\mathrm{ref}}$, a flip persists instead of being corrected on the next cycle.''} Also change ``leaving the primitive's own prior output'' to ``so we use the primitive's own prior output''.

\textbf{M5. The simulated noise model is right, and Eq. (7) is an approximation to it} \\
\textit{Response:} Valid point on spatial correlation and approximations. \\
\textbf{Minimal Fix Proposed:} Add \textit{``gives, approximately, one effective noise level''} before Eq. (7), and explicitly state that $\eta^u$ and $\eta^v$ are independent. Drop the ``validated on the double gyre only'' clause.

\vspace{-0.2cm}
\section*{Further Points (M6 - M9)}

\textbf{M6. Tangent selection off the trench is not analyzed} \\
\textit{Response:} We will adopt the minimal text scope rather than adding the new basin-of-acquisition map. \\
\textbf{Minimal Fix Proposed:} Scope the claim in the Abstract and Section C3 to \textit{``from starts up to $2.7\rho$ off the structure.''}

\textbf{M7. The ocean trial exercises both segment types, but its $s_1$ result depends on the interpolant} \\
\textit{Response:} This is a fair assessment of the bilinear sampler vs. smoother interpolants. \\
\textbf{Minimal Fix Proposed:} State the bilinear sampler in V-A. In V-B, report the eigenvector split and matching references per segment. In Section VI, add the limitation: \textit{``The $s_1$ path depends on the spatial interpolant after about 14 h, while the $D$ path does not.''}

\textbf{M8. Statements true on the benchmark are written as general} \\
\textit{Response:} Agreed. \\
\textbf{Minimal Fix Proposed:} Change ``a flow saddle is...'' to \textit{``on the benchmark, a flow saddle is...''} in III-C, and restrict the strain-dominated claim to \textit{``In incompressible flow, where $D < 0$''} in II-D.

\textbf{M9. The minimality and conic statements need one line each} \\
\textit{Response:} The degenerate conic case needs addressing. \\
\textbf{Minimal Fix Proposed:} Add \textit{``since $\boldsymbol{\Phi} \mathbf{c} = \mathbf{0}$ for nonzero $\mathbf{c}$ exactly when all six satisfy $\mathbf{c}^T\boldsymbol{\phi}(\mathbf{p}) = 0$, degenerate conics such as line pairs included, so any two-row layout is singular''} in II-B.

\vspace{-0.2cm}
\section*{Overclaims and Housekeeping}

\textit{Response:} We agree with all the quick wins and housekeeping points. \\
\textbf{Minimal Fixes Proposed:}
\begin{itemize}
    \setlength{\itemsep}{0pt}
    \setlength{\parskip}{0pt}
    \item Delete the word ``exact'' describing the curves.
    \item Delete the closing flourish ``neither could see''.
    \item Cite or cut the unsupported claim about sampling larger basins without adaptive navigation.
    \item Clarify ambiguous numbers and locations identified in Section 5.2.
    \item Fix the figure order (IEEE chronological requirement), fix the notation collision on $\mathbf{J}$, and update references [2] and [34] as requested.
    \item Remove the inline source comment ``\%\ \#\#\#\#\#\#\ I have proofread...''.
\end{itemize}

\end{document}
